\begin{ouverture}[vecteursplan]{8}{Vecteurs du plan}{Partie 3 · Géométrie}
{Translations · Coordonnées · Sommes · Déterminants}
\objectif{Calculer et représenter les coordonnées d’un vecteur.}
\objectif{Utiliser Chasles et les opérations vectorielles.}
\objectif{Démontrer alignement et parallélisme par un déterminant.}
\prerequis{Repérage du plan et calcul avec les nombres relatifs.}
\end{ouverture}
\pagestyle{otMain}\essentiel\fontsize{11}{15}\selectfont
\essSection{1}{Déplacements et coordonnées}
\begin{definition}Le vecteur $\vect{AB}$, appelé vecteur déplacement de $A$ à $B$, décrit la translation qui envoie $A$ sur $B$. Deux vecteurs égaux ont même direction, même sens et même longueur : $\vect{AB}=\vect{CD}$ signifie que $ABDC$ est un parallélogramme, éventuellement aplati. Toutes les flèches $\vect{AB},\vect{CD},\ldots$ qui représentent un même vecteur $\vec u$ en sont des représentants. Le vecteur nul $\vec0=\vect{AA}$ correspond à un déplacement nul. Si $O$ est l’origine, $\vect{OM}$ est le vecteur position de $M$. Dans un repère orthonormé $(O;\vec i,\vec j)$, la base $(\vec i,\vec j)$ est orthonormée et :
\[\vect{AB}=(x_B-x_A\,;\,y_B-y_A),\qquad\|\vec u\|=\sqrt{u_x^2+u_y^2}.\]
\end{definition}
\begin{propriete}La relation de Chasles est $\vect{AB}+\vect{BC}=\vect{AC}$. Si $\vec u=(a;b)$ et $\vec v=(c;d)$, alors $\vec u+\vec v=(a+c;b+d)$ et $k\vec u=(ka;kb)$.\end{propriete}
\begin{savoirfaire}[Construire une somme]À partir d’une même origine, construire le parallélogramme porté par les deux vecteurs : sa diagonale représente leur somme. On peut aussi reporter le second vecteur à l’extrémité du premier.\end{savoirfaire}
\begin{attention}Les coordonnées d’un point décrivent sa position ; celles d’un vecteur décrivent un déplacement. $\vect{AB}$ et $\vect{BA}$ sont opposés.\end{attention}
\clearpage
\essSection{2}{Colinéarité et choix d’une méthode}
\begin{definition}Deux vecteurs sont colinéaires si l’un est un multiple de l’autre, ou si l’un est nul. Pour $\vec u=(a;b)$ et $\vec v=(c;d)$, leur déterminant est $\det(\vec u,\vec v)=ad-bc$.\end{definition}
\begin{theoreme}Deux vecteurs du plan sont colinéaires si et seulement si leur déterminant est nul. Des points $A,B,C$ sont alignés si et seulement si $\det(\vect{AB},\vect{AC})=0$. Pour deux droites, choisir des vecteurs directeurs non nuls.\end{theoreme}
\begin{propriete}Un point $I$ est le milieu de $[AB]$ si et seulement si $\vect{AI}=\vect{IB}$, ou encore $\vect{IA}+\vect{IB}=\vec0$.\end{propriete}
\begin{savoirfaire}[Dans une base]Si $\vec i=(1;0)$ et $\vec j=(0;1)$, $\vec u=(a;b)$ s’écrit $a\vec i+b\vec j$. Plus généralement, deux vecteurs non colinéaires permettent de décomposer tous les vecteurs du plan.\end{savoirfaire}
\otIntertitre{En physique}
Le vecteur déplacement s’additionne comme les déplacements successifs ; une durée ou une masse s’additionne comme un nombre. Préciser si une donnée est une grandeur vectorielle ou scalaire évite de confondre sens et longueur.
\clearpage\rubrique{Découverte}
\begin{blocLivre}{ex-8-1}
\begin{exercice}[1]{Calculer un déplacement}
On donne $A(1;2)$ et $B(5;-1)$. Calculer $\vect{AB}$ et $\vect{BA}$.
\label{fig-ex-8-1}
\figureGeo{vect-01}
\end{exercice}
\end{blocLivre}
\begin{blocLivre}{ex-8-2}
\begin{exercice}[1]{Additionner des vecteurs}
Pour $\vec u=(3;1)$ et $\vec v=(-1;2)$, calculer et construire $\vec u+\vec v$.
\label{fig-ex-8-2}
\figureGeo{vect-02}
\end{exercice}
\end{blocLivre}
\begin{blocLivre}{ex-8-3}
\begin{exercice}[1]{Multiplier par un réel}
Pour $\vec u=(2;-1)$, calculer $3\vec u$ et $-2\vec u$. Décrire leur sens par rapport à celui de $\vec u$.
\label{fig-ex-8-3}
\figureGeo{vect-03}
\end{exercice}
\end{blocLivre}
\begin{blocLivre}{ex-8-4}
\begin{exercice}[1]{Construire une image par translation}
L’image de $A(-2;1)$ par la translation de vecteur $(4;2)$ est $B$. Déterminer ses coordonnées.
\label{fig-ex-8-4}
\figureGeo{vect-04}
\end{exercice}
\end{blocLivre}
\begin{blocLivre}{ex-8-5}
\begin{exercice}[1]{Une norme}
Calculer la norme de $\vec u=(-3;4)$, puis celle de $2\vec u$.
\label{fig-ex-8-5}
\figureGeo{vect-05}
\end{exercice}
\end{blocLivre}

\clearpage\rubrique{Entraînement}
\begin{blocLivre}{ex-8-6}
\begin{exercice}[2]{La relation de Chasles}
Simplifier $\vect{AB}+\vect{BC}+\vect{CD}$, puis $\vect{AB}+\vect{CA}$.
\label{fig-ex-8-6}
\figureGeo{vect-06}
\end{exercice}
\end{blocLivre}
\begin{blocLivre}{ex-8-7}
\begin{exercice}[2]{Tester la colinéarité}
Les vecteurs $(2;3)$ et $(6;9)$ sont-ils colinéaires ? Même question pour $(2;3)$ et $(4;7)$.
\label{fig-ex-8-7}
\figureGeo{vect-07}
\end{exercice}
\end{blocLivre}
\begin{blocLivre}{ex-8-8}
\begin{exercice}[2]{Trois points alignés ?}
On donne $A(1;1)$, $B(3;4)$ et $C(5;7)$. Démontrer leur alignement et préciser la position de $B$.
\label{fig-ex-8-8}
\figureGeo{vect-08}
\end{exercice}
\end{blocLivre}
\begin{blocLivre}{ex-8-9}
\begin{exercice}[2]{Le quatrième sommet}
$A(0;1)$, $B(4;2)$ et $C(5;5)$ sont trois sommets consécutifs d’un parallélogramme $ABCD$. Déterminer $D$.
\label{fig-ex-8-9}
\figureGeo{vect-09}
\end{exercice}
\end{blocLivre}
\begin{blocLivre}{ex-8-10}
\begin{exercice}[2]{Une combinaison}
Pour $\vec u=(2;1)$ et $\vec v=(-1;3)$, calculer $2\vec u-\vec v$ puis exprimer $(5;-1)$ avec ces vecteurs.
\label{fig-ex-8-10}
\figureGeo{vect-10}
\end{exercice}
\end{blocLivre}

\clearpage\rubrique{Problèmes}
\begin{blocLivre}{ex-8-11}
\begin{exercice}[3]{Un paramètre de colinéarité}
Pour quelle valeur de $m$ les vecteurs $(m;2)$ et $(3;m-1)$ sont-ils colinéaires ? Factoriser l’équation obtenue.
\label{fig-ex-8-11}
\figureGeo{vect-11}
\end{exercice}
\end{blocLivre}
\begin{blocLivre}{ex-8-12}
\begin{exercice}[3]{Un point sur une droite}
$A(1;2)$, $B(5;4)$ et $M(t;5)$. Déterminer $t$ pour que $A,B,M$ soient alignés.
\label{fig-ex-8-12}
\figureGeo{vect-12}
\end{exercice}
\end{blocLivre}
\begin{blocLivre}{ex-8-13}
\begin{exercice}[3]{Une base non usuelle}
On donne $\vec u=(1;1)$ et $\vec v=(1;-1)$. Vérifier qu’ils ne sont pas colinéaires, puis exprimer $\vec w=(4;2)$ comme $a\vec u+b\vec v$.
\label{fig-ex-8-13}
\figureGeo{vect-13}
\end{exercice}
\end{blocLivre}
\begin{blocLivre}{ex-8-14}
\begin{exercice}[3]{Les milieux d’un triangle}
$A(0;0)$, $B(6;2)$ et $C(2;6)$. On donne $I(3;1)$, milieu de $[AB]$, et $J(1;3)$, milieu de $[AC]$. Montrer vectoriellement que $(IJ)$ est parallèle à $(BC)$.
\label{fig-ex-8-14}
\figureGeo{vect-14}
\end{exercice}
\end{blocLivre}
\begin{blocLivre}{ex-8-15}
\begin{exercice}[3]{Un déplacement composé}
Un robot part de $A(1;1)$ et effectue $n$ déplacements $(2;1)$ puis $n-1$ déplacements $(-1;2)$, avec $n\ge1$ entier. Donner l’arrivée $M_n$. Peut-il atteindre $(8;8)$ ?
\label{fig-ex-8-15}
\figureGeo{vect-15}
\end{exercice}
\end{blocLivre}

\clearpage
